% Work-MAT | 2. ročník | hodina 42 | 2026-11-22
% Zadání vybírána z WorkMAT_Logaritmy_Sbirka_PodleMetod.tex (481 příkladů)
% Vždy zachována matematická formulace zdroje, učitelův výklad a řešení dopsány k vybraným příkladům.
\documentclass[11pt,a4paper]{article}
\input{../../sablony/pripravy/priprava.sty}
\usepackage{mathtools}
\nastavPripravu{2. ročník}{Logaritmické nerovnice}{42}{Základní logaritmické nerovnice}
\begin{document}
\blok{Návaznost a cíl hodiny}
Od rovnic přecházíme k nerovnicím. Žáci už znají, že logaritmická funkce může být rostoucí nebo klesající; rozhodující je přípustný základ a definiční obor argumentu.
\poznamka{\textbf{Časový plán (45--55 minut):} Úvod 5 min, příklady 1--2 17 min, příklad 3 8 min, samostatná práce a--e 13 min, závěr 5 min. Příklad 4 dle rychlosti.}

\blok{Klíčová metoda}
\[
\begin{cases}a>1: & \log_a u<\log_a v\iff u<v,\\0<a<1: & \log_a u<\log_a v\iff u>v,\end{cases}\quad u,v>0.
\]
\poznamka{Nejde o novou definici. Jde o důsledek prostoty a pravidel pro logaritmy, která už žáci znají. V řešení mají vždy uvést podmínky.}

\blok{Společné řešené příklady}

% Sbírka: oddíl 18, položka 9
\needspace{9\baselineskip}
\begin{prikladReseniBox}[title={Příklad 1}]
Vyřeš: $\displaystyle \log_2(x+1)<2$
\tcblower
\[
\begin{aligned}
&x+1>0\Longrightarrow x>-1,\\
\log_2(x+1)&<\log_2 4,\\
x+1&<4,\\
K&=(-1;3).
\end{aligned}
\]
\end{prikladReseniBox}
\poznamka{Z podmínky a z nerovnice se bere průnik, ne pouze vyřešená algebraická část.}
% Sbírka: oddíl 18, položka 10
\needspace{9\baselineskip}
\begin{prikladReseniBox}[title={Příklad 2}]
Vyřeš: $\displaystyle \log_{\frac12}(2x-1)<-1$
\tcblower
\[
\begin{aligned}
&2x-1>0,\\
\log_{\frac12}(2x-1)&<\log_{\frac12}2,\\
2x-1&>2,\\
K&=\left(\frac32;\infty\right).
\end{aligned}
\]
\end{prikladReseniBox}
\poznamka{Při základu mezi nulou a jednou se znaménko při odlogaritmování obrací.}
% Sbírka: oddíl 18, položka 11
\needspace{9\baselineskip}
\begin{prikladReseniBox}[title={Příklad 3}]
Vyřeš: $\displaystyle \log_3(x+2)<\log_3(3-x)$
\tcblower
\[
\begin{aligned}
&-2<x<3,\\
\log_3(x+2)&<\log_3(3-x),\\
x+2&<3-x,\\
x&<\frac12,\\
K&=\left(-2;\frac12\right).
\end{aligned}
\]
\end{prikladReseniBox}
\poznamka{Dvě současné podmínky zapisuj na číselnou osu. V součtu vždy najdi správný průnik.}
% Sbírka: oddíl 19, položka 13
\needspace{9\baselineskip}
\begin{prikladReseniBox}[title={Příklad 4}]
Vyřeš: $\displaystyle \log_{0{,}5}x+\log_{0{,}5}(x+3)\ge2\log_{0{,}5}2$
\tcblower
\[
\begin{aligned}
&x>0,\\
\log_{\frac12}\bigl(x(x+3)\bigr)&\ge\log_{\frac12}4,\\
x^2+3x&\le4,\\
(x+4)(x-1)&\le0,\\
K&=(0;1\rangle.
\end{aligned}
\]
\end{prikladReseniBox}
\poznamka{Před použitím vzorce pro součin logaritmů se musí ověřit kladnost obou původních argumentů. Intervaly zapisuj českou notací.}
\blok{Samostatná práce}
\poznamka{Ve třídě doporučuji nejdříve úlohy a)--d); další jsou rezervou pro rychlejší žáky nebo volitelným dokončením po hodině. Neopisovat řešený vzor — vybrat postup.}
\begin{zadaniBox}[title={Samostatná práce 1}]
\begin{enumerate}[label=\alph*)]
% Sbírka: oddíl 18, položka 1
\item $\displaystyle \log_3(x-4)<0$
% Sbírka: oddíl 18, položka 2
\item $\displaystyle \log_5(x^2-2x+1)\ge0$
% Sbírka: oddíl 18, položka 3
\item $\displaystyle \log_{\frac12}(x+5)\le0$
% Sbírka: oddíl 18, položka 4
\item $\displaystyle \log_{\frac14}(4x^2+3x)>0$
% Sbírka: oddíl 18, položka 6
\item $\displaystyle \log_5(2x-1)>-1$
% Sbírka: oddíl 18, položka 7
\item $\displaystyle \log_{\frac12}(2x+4)\ge-3$
% Sbírka: oddíl 18, položka 12
\item $\displaystyle \log_3(x-1)>\log_3 3$
% Sbírka: oddíl 18, položka 13
\item $\displaystyle \log_5(2x+3)\le1$
\end{enumerate}
\end{zadaniBox}
\blok{Úkol na doma}
Vyřeš úlohy a)--e) a zapiš podmínky. Zbývající úlohy jsou procvičování navíc; podle rychlosti je můžeš dokončit doma. 
\begin{zadaniBox}[title={Domácí úkol}]
\begin{enumerate}[label=\alph*)]
% Sbírka: oddíl 18, položka 5
\item $\displaystyle \log_3(x+4)\le4$
% Sbírka: oddíl 18, položka 8
\item $\displaystyle \log_{0{,}1}(x-6)<2$
% Sbírka: oddíl 18, položka 14
\item $\displaystyle \log_{10}(x-3)+\log_{10}(2x-1)<0$
% Sbírka: oddíl 18, položka 15
\item $\displaystyle \log_2(x+1)\ge\log_2(5-x)$
% Sbírka: oddíl 18, položka 16
\item $\displaystyle 2\log_4(x-1)-\log_4(x+3)>1$
% Sbírka: oddíl 18, položka 17
\item $\displaystyle \log_{0{,}5}(x-2)+\log_{0{,}5}(x-4)\le-2$
% Sbírka: oddíl 18, položka 18
\item $\displaystyle \log_3(x+2)<\log_3(7-x)$
% Sbírka: oddíl 18, položka 19
\item $\displaystyle 3\log_2(x-1)+\log_2(x+3)\ge5$
\end{enumerate}
\end{zadaniBox}
\poznamka{\textbf{Učitelovo vodítko:} U každé nerovnice nejdříve napiš definiční obor. Při základu mezi nulou a jednou měň směr nerovnosti a nakonec vezmi průnik. Domácí sada má 8 různých zadání; základních pět úloh na přibližně 15--30 minut a zbytek jako volitelnou zásobu.}
\end{document}
